% ---
% titre: Recette de crêpes (tableau de proportionnalité, quatrième proportionnelle)
% theme: FA
% chapitre: Proportionnalité
% sous_chapitre: Résoudre des problèmes de proportionnalité
% niveau: [10e]
% type: EVAL
% tags: [c-proportionnalite,p-question-TS]
% difficulte: 1
% corrige: true
% ---

\question
Voici la recette de crêpes de Julie pour 4 personnes.

\vspace{5pt}\def\arraystretch{2}
\begin{tabularx}{\linewidth}{|p{3cm}|C|}
\hline
\rowcolor{gray!20}\textbf{Ingrédient} & \textbf{Pour 4 personnes}\\
\hline
Farine & 300\,g\\
\hline
Lait & 0,6\,l\\
\hline
Oeufs & 2\\
\hline
Sucre & 60\,g\\
\hline
\end{tabularx}
\def\arraystretch{1}

\begin{parts}
\vspace{20pt}\part[2] Combien de sucre faut-il pour 10 personnes?

\vspace{10pt}{\footnotesize\raggedleft Espace pour ta démarche et tes calculs\par}
\begin{solutionorgrid}[120pt]

\vspace{5pt}
\def\arraystretch{1.5}
\begin{tabularx}{0.5\textwidth}{l|C|C|}
Sucre [g]&60&\\
\hline
Nombre de personnes&4&10\\
\end{tabularx}
\def\arraystretch{1}

\hfill\textbf{(0,5pt)} pour la construction du tableau

\vspace{15pt}$10 \cdot 60 : 4 = 150$\,g \hfill calcul \textbf{(0,5\,pt)} et réponse \textbf{(0,5\,pt)}
\end{solutionorgrid}

\begin{tcolorbox}[enhanced jigsaw, interior hidden, colframe=box, parbox=false]%
Réponse: 
\sol{\vspace{20pt}}{Pour 10 personnes, il faut 150\,g de sucre. \textbf{(0,5\,pt)}}
\end{tcolorbox}

\vspace{20pt}\part[2] Julie a 1300\,g de farine. Pour combien de personnes peut-elle, au maximum, préparer des crêpes avec cette recette ?

\vspace{10pt}{\footnotesize\raggedleft Espace pour ta démarche et tes calculs\par}
\begin{solutionorgrid}[120pt]

\vspace{5pt}
\def\arraystretch{1.5}
\begin{tabularx}{0.5\textwidth}{l|C|C|}
Farine [g]&300&1300\\
\hline
Nombre de personnes&4&\\
\end{tabularx}
\def\arraystretch{1}

\hfill\textbf{(0,5pt)} pour la construction du tableau

\vspace{15pt}$4 \cdot 1300 : 300 \approx 17,3$\,personne \hfill calcul \textbf{(0,5\,pt)} et réponse \textbf{(0,5\,pt)}
\end{solutionorgrid}

\begin{tcolorbox}[enhanced jigsaw, interior hidden, colframe=box, parbox=false]%
Réponse: 
\sol{\vspace{20pt}}{Julie peut préparer des crêpes pour 17 personnes au maximum. \textbf{(0,5\,pt)}}
\end{tcolorbox}

\vspace{20pt}\part[2] Julie a aussi 1 litre de lait. Est-ce suffisant pour le nombre de personnes trouvé en (b)? Justifie ta réponse.

{\footnotesize Si tu n’as pas trouvé la réponse à la question précédente, utilise 15 personnes.}

\vspace{10pt}{\footnotesize\raggedleft Espace pour ta démarche et tes calculs\par}
\begin{solutionorgrid}[120pt]

\vspace{5pt}
\def\arraystretch{1.5}
\begin{tabularx}{0.5\textwidth}{l|C|C|}
Lait [l]&0,6&\\
\hline
Nombre de personnes&4&17\\
\end{tabularx}
\def\arraystretch{1}

\hfill\textbf{(0,5pt)} pour la construction du tableau

\vspace{15pt}$17 \cdot 0,6 : 4 = 2,55$\,litres \hfill calcul \textbf{(0,5\,pt)} et réponse \textbf{(0,5\,pt)}
\end{solutionorgrid}

\begin{tcolorbox}[enhanced jigsaw, interior hidden, colframe=box, parbox=false]%
Réponse: 
\sol{\vspace{20pt}}{Non, ce n'est pas suffisant : il faut 2,55\,l de lait et elle n'en a que 1\,l. \textbf{(0,5\,pt)}}
\end{tcolorbox}
\end{parts}
